\documentclass[10pt,a4paper]{amsart}
\usepackage{amsmath,amssymb,amsthm,graphicx}
\usepackage{verbatim}
\newenvironment{reference}{\comment}{\endcomment}
%\newenvironment{reference}{}{}
\newenvironment{slogan}{\comment}{\endcomment}
\newenvironment{history}{\comment}{\endcomment}

% For commutative diagrams we use Xy-pic
\usepackage[all]{xy}

% We use 2cell for 2-commutative diagrams.
\xyoption{2cell}
% Theorem environments.
%
\theoremstyle{plain}
\newtheorem{theorem}[subsection]{Theorem}
\newtheorem{proposition}[subsection]{Proposition}
\newtheorem{lemma}[subsection]{Lemma}

\theoremstyle{definition}
\newtheorem{definition}[subsection]{Definition}
\newtheorem{example}[subsection]{Example}
\newtheorem{exercise}[subsection]{Exercise}
\newtheorem{situation}[subsection]{Situation}

\theoremstyle{remark}
\newtheorem{remark}[subsection]{Remark}
\newtheorem{remarks}[subsection]{Remarks}

\numberwithin{equation}{subsection}

% Macros
%
\def\lim{\mathop{\mathrm{lim}}\nolimits}
\def\colim{\mathop{\mathrm{colim}}\nolimits}
\def\Spec{\mathop{\mathrm{Spec}}}
\def\Hom{\mathop{\mathrm{Hom}}\nolimits}
\def\Ext{\mathop{\mathrm{Ext}}\nolimits}
\def\SheafHom{\mathop{\mathcal{H}\!\mathit{om}}\nolimits}
\def\SheafExt{\mathop{\mathcal{E}\!\mathit{xt}}\nolimits}
\def\Sch{\mathit{Sch}}
\def\Mor{\mathop{\mathrm{Mor}}\nolimits}
\def\Ob{\mathop{\mathrm{Ob}}\nolimits}
\def\Sh{\mathop{\mathit{Sh}}\nolimits}
\def\NL{\mathop{N\!L}\nolimits}
\def\CH{\mathop{\mathrm{CH}}\nolimits}
\def\proetale{{pro\text{-}\acute{e}tale}}
\def\etale{{\acute{e}tale}}
\def\QCoh{\mathit{QCoh}}
\def\Ker{\mathop{\mathrm{Ker}}}
\def\Im{\mathop{\mathrm{Im}}}
\def\Coker{\mathop{\mathrm{Coker}}}
\def\Coim{\mathop{\mathrm{Coim}}}

% Boxtimes
%
\DeclareMathSymbol{\boxtimes}{\mathbin}{AMSa}{"02}

%
% Macros for moduli stacks/spaces
%
\def\QCohstack{\mathcal{QC}\!\mathit{oh}}
\def\Cohstack{\mathcal{C}\!\mathit{oh}}
\def\Spacesstack{\mathcal{S}\!\mathit{paces}}
\def\Quotfunctor{\mathrm{Quot}}
\def\Hilbfunctor{\mathrm{Hilb}}
\def\Curvesstack{\mathcal{C}\!\mathit{urves}}
\def\Polarizedstack{\mathcal{P}\!\mathit{olarized}}
\def\Complexesstack{\mathcal{C}\!\mathit{omplexes}}
% \Pic is the operator that assigns to X its picard group, usage \Pic(X)
% \Picardstack_{X/B} denotes the Picard stack of X over B
% \Picardfunctor_{X/B} denotes the Picard functor of X over B
\def\Pic{\mathop{\mathrm{Pic}}\nolimits}
\def\Picardstack{\mathcal{P}\!\mathit{ic}}
\def\Picardfunctor{\mathrm{Pic}}
\def\Deformationcategory{\mathcal{D}\!\mathit{ef}}
\usepackage{cleveref}
\crefname{originalproposition}{original proposition}{original propositions}
\Crefname{originalproposition}{Original proposition}{Original propositions}
\makeatletter
\newlabel{proposition-finite-flat-equivalence-global}{{07S6}{source}{}{sourcepointer-1}{}}
\newlabel{proposition-finite-flat-equivalence-global@cref}{{[originalproposition][0][]07S6}{[1][0][]source}{}{}{}}
\makeatother
\crefname{originallemma}{original lemma}{original lemmas}
\Crefname{originallemma}{Original lemma}{Original lemmas}
\makeatletter
\newlabel{lemma-quotient-scheme}{{07S7}{source}{}{sourcepointer-2}{}}
\newlabel{lemma-quotient-scheme@cref}{{[originallemma][0][]07S7}{[1][0][]source}{}{}{}}
\makeatother
\crefname{originalreference}{original reference}{original references}
\Crefname{originalreference}{Original reference}{Original references}
\makeatletter
\newlabel{descent-lemma-descending-property-quasi-compact}{{02KQ}{source}{}{sourcepointer-3}{}}
\newlabel{descent-lemma-descending-property-quasi-compact@cref}{{[originalreference][0][]02KQ}{[1][0][]source}{}{}{}}
\makeatother
\crefname{originalreference}{original reference}{original references}
\Crefname{originalreference}{Original reference}{Original references}
\makeatletter
\newlabel{descent-lemma-descending-property-closed-immersion}{{02L6}{source}{}{sourcepointer-4}{}}
\newlabel{descent-lemma-descending-property-closed-immersion@cref}{{[originalreference][0][]02L6}{[1][0][]source}{}{}{}}
\makeatother
\crefname{originalreference}{original reference}{original references}
\Crefname{originalreference}{Original reference}{Original references}
\makeatletter
\newlabel{schemes-lemma-quasi-compact-permanence}{{03GI}{source}{}{sourcepointer-5}{}}
\newlabel{schemes-lemma-quasi-compact-permanence@cref}{{[originalreference][0][]03GI}{[1][0][]source}{}{}{}}
\makeatother
\crefname{originalreference}{original reference}{original references}
\Crefname{originalreference}{Original reference}{Original references}
\makeatletter
\newlabel{etale-lemma-characterize-closed-immersion}{{04XV}{source}{}{sourcepointer-6}{}}
\newlabel{etale-lemma-characterize-closed-immersion@cref}{{[originalreference][0][]04XV}{[1][0][]source}{}{}{}}
\makeatother
\crefname{originalreference}{original reference}{original references}
\Crefname{originalreference}{Original reference}{Original references}
\makeatletter
\newlabel{morphisms-lemma-image-proper-scheme-closed}{{01W6}{source}{}{sourcepointer-7}{}}
\newlabel{morphisms-lemma-image-proper-scheme-closed@cref}{{[originalreference][0][]01W6}{[1][0][]source}{}{}{}}
\makeatother
\crefname{originaldefinition}{original definition}{original definitions}
\Crefname{originaldefinition}{Original definition}{Original definitions}
\makeatletter
\newlabel{definition-separated}{{03BS}{source}{}{sourcepointer-8}{}}
\newlabel{definition-separated@cref}{{[originaldefinition][0][]03BS}{[1][0][]source}{}{}{}}
\makeatother
\title{Separatedness properties descend along finite locally free groupoid quotients}
\author{The Stacks Project Authors}
\date{stacks-project:spaces-properties.tex:local-snapshot-20261012; GNU FDL 1.2 or later}
\begin{document}
\maketitle
\noindent\textit{Editorial note.} All delivered mathematical text is the author's own, copied verbatim from the pinned Stacks Project chapter file spaces-properties.tex; no statement or proof step was rewritten, reordered or removed.\par\medskip
\noindent\textit{Editorial note.} Span recomputation by label pairing (LANE\_BRIEF 10.2): in this chapter the author's proof environment follows the statement environment as a standalone block; the spec covers the complete environment pair (statement span, proof span) as recorded in the pool candidate.\par\medskip
\noindent\textit{Editorial note.} Occupation scan: combined claim\_key + statement\_anchor + generic-line scan over all registration files (active and withdrawn) found 13 S10-geom published cards for spaces-properties.tex (prop-occupied.json, 2026-10-14); those 13 labels are excluded from this batch by design. For the 77 free labels: 0 occupied. The registry mentions field (7557 files) and focused-seeds (645) both point to the same S10-geom cards (dual-source check), no other occupation.\par\medskip
\ifdefined\hypertarget\hypertarget{sourcepointer-1}{}\fi
\noindent\textit{Source reference.} Original proposition 07S6: sources/stacks-chow/tags.txt, source lines 9891--9891. Cross-chapter standard prerequisite cited by the author inside the counted span; it is used as a pointer and its own proof is not part of this item.\par\medskip
\ifdefined\hypertarget\hypertarget{sourcepointer-2}{}\fi
\noindent\textit{Source reference.} Original lemma 07S7: sources/stacks-chow/tags.txt, source lines 9892--9892. Cross-chapter standard prerequisite cited by the author inside the counted span; it is used as a pointer and its own proof is not part of this item.\par\medskip
\ifdefined\hypertarget\hypertarget{sourcepointer-3}{}\fi
\noindent\textit{Source reference.} Original reference 02KQ: sources/stacks-chow/tags.txt, source lines 3381--3381. Cross-chapter standard prerequisite cited by the author inside the counted span; it is used as a pointer and its own proof is not part of this item.\par\medskip
\ifdefined\hypertarget\hypertarget{sourcepointer-4}{}\fi
\noindent\textit{Source reference.} Original reference 02L6: sources/stacks-chow/tags.txt, source lines 3397--3397. Cross-chapter standard prerequisite cited by the author inside the counted span; it is used as a pointer and its own proof is not part of this item.\par\medskip
\ifdefined\hypertarget\hypertarget{sourcepointer-5}{}\fi
\noindent\textit{Source reference.} Original reference 03GI: sources/stacks-chow/tags.txt, source lines 4529--4529. Cross-chapter standard prerequisite cited by the author inside the counted span; it is used as a pointer and its own proof is not part of this item.\par\medskip
\ifdefined\hypertarget\hypertarget{sourcepointer-6}{}\fi
\noindent\textit{Source reference.} Original reference 04XV: sources/stacks-chow/tags.txt, source lines 6408--6408. Cross-chapter standard prerequisite cited by the author inside the counted span; it is used as a pointer and its own proof is not part of this item.\par\medskip
\ifdefined\hypertarget\hypertarget{sourcepointer-7}{}\fi
\noindent\textit{Source reference.} Original reference 01W6: sources/stacks-chow/tags.txt, source lines 2471--2471. Cross-chapter standard prerequisite cited by the author inside the counted span; it is used as a pointer and its own proof is not part of this item.\par\medskip
\ifdefined\hypertarget\hypertarget{sourcepointer-8}{}\fi
\noindent\textit{Source reference.} Original definition 03BS: sources/stacks-chow/tags.txt, source lines 4359--4359. Cross-chapter standard prerequisite cited by the author inside the counted span; it is used as a pointer and its own proof is not part of this item.\par\medskip
\begin{lemma}
\label{lemma-quotient-separated}
Notation and assumptions as in
Proposition \ref{proposition-finite-flat-equivalence-global}. Then
\begin{enumerate}
\item if $U$ is quasi-separated over $S$, then $U/R$ is quasi-separated
over $S$,
\item if $U$ is quasi-separated, then $U/R$ is quasi-separated,
\item if $U$ is separated over $S$, then $U/R$ is separated over $S$,
\item if $U$ is separated, then $U/R$ is separated, and
\item add more here.
\end{enumerate}
Similar results hold in the setting of Lemma \ref{lemma-quotient-scheme}.
\end{lemma}

\begin{proof}
Since $M$ represents the quotient sheaf we have a cartesian diagram
$$
\xymatrix{
R \ar[r]_-j \ar[d] & U \times_S U \ar[d] \\
M \ar[r] & M \times_S M
}
$$
of schemes. Since $U \times_S U \to M \times_S M$ is surjective finite locally
free, to show that $M \to M \times_S M$ is quasi-compact, resp.\ a closed
immersion, it suffices to show that $j : R \to U \times_S U$ is
quasi-compact, resp.\ a closed immersion, see
Descent, Lemmas \ref{descent-lemma-descending-property-quasi-compact} and
\ref{descent-lemma-descending-property-closed-immersion}.
Since $j : R \to U \times_S U$ is a morphism over $U$ and since
$R$ is finite over $U$, we see that $j$ is quasi-compact as soon
as the projection $U \times_S U \to U$ is quasi-separated
(Schemes, Lemma \ref{schemes-lemma-quasi-compact-permanence}).
Since $j$ is a monomorphism and locally of finite type, we see that
$j$ is a closed immersion as soon as it is proper
(\'Etale Morphisms, Lemma \ref{etale-lemma-characterize-closed-immersion})
which will be the case as soon as the projection
$U \times_S U \to U$ is separated
(Morphisms, Lemma \ref{morphisms-lemma-image-proper-scheme-closed}).
This proves (1) and (3). To prove (2) and (4) we replace $S$ by
$\Spec(\mathbf{Z})$, see Definition \ref{definition-separated}.
Since Lemma \ref{lemma-quotient-scheme} is proved through an application of
Proposition \ref{proposition-finite-flat-equivalence-global}
the final statement is clear too.
\end{proof}

\end{document}
